\documentclass[11pt]{article}

\usepackage[T1]{fontenc}
\usepackage{lmodern}
\usepackage{microtype}
\usepackage[margin=1.05in]{geometry}
\usepackage{amsmath,amssymb,amsthm,mathtools}
\usepackage[shortlabels]{enumitem}
\usepackage{booktabs}
\usepackage{xcolor}
\usepackage[colorlinks=true,linkcolor=blue!45!black,
  citecolor=blue!45!black,urlcolor=blue!45!black]{hyperref}
\hypersetup{
  pdftitle={New Foundations Refutes the Weak Partition Principle},
  pdfauthor={},
  pdfsubject={Weak choice, Hartogs numbers, and finite typed blocks in New Foundations},
  pdfkeywords={New Foundations, weak partition principle, Hartogs number,
    first differences, finite typed replay}
}
\setlength{\emergencystretch}{2em}

\newtheorem{theorem}{Theorem}[section]
\newtheorem{proposition}[theorem]{Proposition}
\newtheorem{lemma}[theorem]{Lemma}
\newtheorem{corollary}[theorem]{Corollary}
\numberwithin{equation}{section}
\theoremstyle{definition}
\newtheorem{definition}[theorem]{Definition}
\theoremstyle{remark}
\newtheorem{remark}[theorem]{Remark}

\newcommand{\NF}{\mathsf{NF}}
\newcommand{\WPP}{\mathsf{WPP}}
\newcommand{\PP}{\mathsf{PP}}
\newcommand{\Pow}{\mathcal P}
\newcommand{\Ord}{\mathrm{Ord}}
\newcommand{\inj}{\hookrightarrow}
\newcommand{\surj}{\twoheadrightarrow}
\newcommand{\otp}{\operatorname{otp}}
\newcommand{\im}{\operatorname{im}}

\title{New Foundations Refutes\\the Weak Partition Principle}
\author{}
\date{August 2026}

\begin{document}
\maketitle

\begin{abstract}
We prove that New Foundations refutes the weak partition principle.  The
proof uses one carrier-coded first-difference construction, one explicitly
shifted WPP padding map, and one canonical comparison of Hartogs quotient
codes in a common ambient base.  This supplies the homogeneous inequality
needed by a finite-orbit contradiction.  Every type displacement is shown:
the source and target of the decisive injection are raised to the same
level before any iterate of the cardinal type operator is cancelled.
Pointwise WPP witnesses are used only to prove image coverage; the final
map is a witness-independent quotient graph, so no choice of a family of
witnesses occurs.  The argument is a direct finite typed replay in NF.  It
uses no direct-limit model, satisfaction predicate, interpretation of
bounded Zermelo set theory, forcing, Replacement, or Axiom of Counting.
\end{abstract}

\section{Introduction}

The weak partition principle is the assertion that a surjection does not
exist in the direction opposite to a strict cardinal inequality.  We use
the following local form:
\begin{equation}
 X\leq^*Y\ \wedge\ Y\leq X
 \quad\Longrightarrow\quad X\leq Y,
 \label{eq:wpp-local}
\end{equation}
where $X\leq Y$ means that $X$ injects into $Y$, and $X\leq^*Y$
means that $X=\varnothing$ or $Y$ surjects onto $X$.

\begin{theorem}[Weak choice fails in NF]\label{thm:main}
\[
 \NF\vdash\neg\WPP.
\]
\end{theorem}

The proof has three parts.  First, assume for contradiction that
$\aleph(\Pow^2\aleph(X))$ is a proper initial segment of
$\aleph(\Pow X)$.  Apply first differences to that one initial segment.
Weak choice embeds the resulting well-ordered first-difference carrier
into $X^2\times\omega$; two direct-image maps then embed the original
initial segment into $\Pow^2\aleph(X)$, contradicting the defining
property of $\aleph(\Pow^2\aleph(X))$.  Second, an explicit finite typed replay turns
the pointwise constructions into one witness-independent Hartogs quotient
graph.  The exact homogeneous inequality is
\begin{equation}
 h(ea)\leq h\bigl(e^2(ha)\bigr).
 \label{eq:intro-homogeneous}
\end{equation}
Finally, specializing $a$ along the universal $T$-chain produces a
cardinal $\alpha$ with
\[
 T\alpha<\alpha<F(T\alpha),
 \qquad F(\beta)=h\bigl(h(e^2\beta)\bigr),
\]
and a least-orbit argument gives an impossible equation on a natural
number.

The finite-block step is included explicitly because it is where an
informal phrase such as ``after the stratification dust settles'' can hide
an invalid one-sided cancellation of $T$.  Our translation raises the two
endpoints to one common level before lowering the injection graph.  Every
later cancellation therefore removes the same iterate of $T$ from both
sides.

\section{Hartogs numbers and the type operator}

For a set $X$, define its Hartogs number by
\[
 \aleph(X)=\min\{\alpha\in\Ord:\alpha\not\leq X\}.
\]
It is obtained from the quotient set of well-orders on subsets of $X$ by
isomorphism, in a fixed finite power-set neighborhood of $X$.  Although
Lindenbaum numbers are useful in a more global presentation of the
argument, the one-witness proof below never forms one.

Let $T$ be the usual type-raising operation on NF-cardinals.  If
$a=|A|$, let $e(a)$ and $h(a)$ be the homogeneous exponential and
Hartogs operations, normalized by
\begin{equation}
 |\Pow A|=T(ea),
 \qquad
 |\mathbf H(A)|=T^2(ha),
 \label{eq:normalization}
\end{equation}
where $\mathbf H(A)$ is the standard typed Hartogs quotient object.
The operations $e$ and $h$ commute with $T$ wherever defined.

We use the following standard facts of the Specker cardinal calculus
\cite{Specker,Forster,HolmesBook}.

\begin{lemma}[Basic $T$-calculus]\label{lem:t-calculus}
In NF:
\begin{enumerate}[(i)]
\item $T$ is an order embedding on cardinals.  On well-orderable
      cardinals, the range of each fixed finite iterate $T^r$ is initial.
\item The restriction of $T$ to the natural numbers is onto, preserves
      order and arithmetic, and preserves parity.
\item If $G$ is a homogeneous partial operation, finite iteration is
      homogeneous:
      \[
       T(G^m(\beta))=G^{Tm}(T\beta)
      \]
      whenever either side is defined.
\item The double-singleton comparison gives $\beta\leq e^2\beta$ on
      well-orderable cardinals in the common domain.
\end{enumerate}
\end{lemma}

\section{Finite typed supports and local maps}

This section gives the finite typed bookkeeping used in the proof.  We do
not construct a one-sorted limit model and do not define a satisfaction
relation.  Every set below is obtained by a displayed stratified
comprehension in NF.

Two pair codes must be distinguished.  Write
$\langle x,y\rangle_0$ for the standard type-level Quine--Rosser pair used
in function graphs, and $\langle\!\langle x,y\rangle\!\rangle$ for the
Kuratowski pair.  The former preserves the common type of its coordinates;
the latter raises it by two.  Write $A\times_KB$ for the Kuratowski product.
The type-level pair is a conservative definitional extension of NF
\cite{Forster,HolmesBook}.  Keeping these pair codes separate is what makes
the shift ledger below exact.

For that diagram, the adjacent type-copy operations are defined externally
and only on the displayed supports.  Put $s(x)=\{x\}$ and recursively
define $s^{r+1}(x)=\{s^r(x)\}$.  For a displayed carrier $A$, put
\[
 s^rA=\{s^r(x):x\in A\}.
\]
Thus $x\in A$ iff $s^r(x)\in s^rA$.  Each image of a displayed set
exists by one stratified comprehension instance; no global graph
$x\mapsto s^r(x)$ is asserted to be a set.  For two coordinate levels the
pairing transport is the metasyntactic map
\[
 \langle x,y\rangle\longmapsto
 \langle s^i(x),s^j(y)\rangle,
\]
expanded into the fixed pair code at the assigned types.  Its inverse on
the supported product is coordinatewise and unique.

Two qualifications are essential.  First, products are transported by a
typed pairing isomorphism: for Kuratowski pairs one does \emph{not} have
$\langle\{x\},\{y\}\rangle=\{\langle x,y\rangle\}$.  Second, forward
naturality of well-order codes is not enough; every supported raised code
and every isomorphism witness used below must lower as well.

\begin{lemma}[Supported local lowering]
\label{lem:local-lowering}
At any fixed compatible assignment of types, the following hold.
\begin{enumerate}[(i)]
\item A subset of an elementwise-shifted displayed carrier lowers uniquely
      by bounded comprehension.
\item A relation between two possibly differently shifted carriers lowers
      uniquely after conjugation by the coordinatewise typed pairing
      isomorphism.  Function, injection, surjection, and bijection are
      preserved and reflected.
\item A well-order on a supported carrier lowers to a well-order; an
      isomorphism between two supported well-order codes lowers; and the
      induced map on the corresponding Hartogs quotient classes is onto as
      well as one-to-one and order preserving.
\end{enumerate}
The assertions remain true after any further common shift.
\end{lemma}

\begin{proof}
For (i), a subset $B$ of the raised support of $X$ lowers to
\[
 \{x\in X:s(x)\in B\}.
\]
At the fixed assigned levels this is stratified comprehension, and raising
the result recovers $B$.  For (ii), transport the pair code by
$\langle x,y\rangle\mapsto\langle s_i(x),s_j(y)\rangle$ and apply (i) to
the supported product.  The four graph properties then transfer
pointwise.

For (iii), every subset of the raised well-order carrier lowers by (i).
Thus a subset witnessing failure of well-foundedness would lower, and
well-foundedness is reflected.  Relation codes and isomorphism graphs lower
by (ii).  Finally every well-order code on a supported raised subset lowers,
so the map on isomorphism classes has reverse coverage, not merely forward
naturality.  Strict-segment order is defined by the same supported graph
formulas and is therefore preserved and reflected.
\end{proof}

Iterating part (iii) gives a natural order isomorphism
\begin{equation}
 \theta_C^r:s^r\mathbf H(C)\longrightarrow\mathbf H(s^rC).
 \label{eq:hartogs-support-naturality}
\end{equation}
If $A\subseteq C$, retaining a well-order code while enlarging its ambient
base gives a set-coded injection
\begin{equation}
 \iota_A^C:\mathbf H(A)\longrightarrow\mathbf H(C).
 \label{eq:hartogs-inclusion}
\end{equation}
Concretely, its graph consists of the pairs
$\langle[u]_A,[u]_C\rangle_0$.  It is well defined and injective because
isomorphism of codes depends only on their carriers and relations, not on
the unused part of the ambient base.

\begin{lemma}[Common-ambient quotient graph]
\label{lem:common-ambient-quotient-graph}
Let $A,B\subseteq C$ be displayed carriers at one level.  If
\[
 \iota_A^C\,``\,\mathbf H(A)
 \subseteq
 \iota_B^C\,``\,\mathbf H(B),
\]
then
\[
 G_{A,B;C}=(\iota_B^C)^{-1}\circ\iota_A^C
\]
is a set-coded injection $\mathbf H(A)\inj\mathbf H(B)$.
\end{lemma}

\begin{proof}
The range inclusion makes the composite total.  Both inclusion maps are
one-to-one, so the inverse relation is functional on the relevant range
and the composite is one-to-one.  The graph mentions no witness used to
prove the range inclusion.
\end{proof}

The complete carrier ledger is as follows.  Type-level function graphs do
not change the carrier level of their aligned endpoints; the occurrences
of $\times_K$ are exactly the places where two levels are added.
\begin{table}[ht]
\centering
\caption{Occurrence-level certificate for the finite replay.}
\label{tab:intrinsic-certificate}
\begin{tabular}{@{}cl@{}}
\toprule
relative level & displayed carriers \\
\midrule
$0$ & $X,\ \Omega$ \\
$1$ & $\Pow X,\ I,\ r_x,\ D,\ q_x,\ R,\ R\mathbin{\upharpoonright}D$ \\
$2$ & $U=X\times_KX,\ V=sD,\ \mathbf H(X),\ \Pow D,\ C_A,\ s^2\Omega$ \\
$3$ & $\Pow^2D,\ \mathbf H(\Pow X),\ s^2I$ \\
$4$ & $P=U\times_Ks^2\Omega,\ s^3D,\ s^3I,\ s(\Pow^2D),$
       $\Pow^2(sD),\ \Pow^2\mathbf H(X)$ \\
$6$ & $s^2P,\ Q,\ \mathbf H(P),\ s^4\mathbf H(X),\ s^5D,$
       $s^3\mathbf H(\Pow X)$ \\
\bottomrule
\end{tabular}
\end{table}
Adding a common external offset makes all levels nonnegative.  The row,
first-difference, column, parity, tagging, and least-preimage formulas have
the displayed stratifications; thus each set in the table is an ordinary
NF comprehension instance.

\begin{lemma}[Typed WPP padding]
\label{lem:typed-wpp-padding}
Let $g:U\surj V$, where $U$ and $V$ have a common carrier level $d$, and
let $\Omega_d$ be the Peano carrier at that level.  Under WPP there is an
injection
\begin{equation}
 s^2V\inj U\times_K\Omega_d.
 \label{eq:typed-padding}
\end{equation}
The conclusion is also valid when $V=\varnothing$.
\end{lemma}

\begin{proof}
The empty case is immediate.  Put
\[
 P=U\times_K\Omega_d,\qquad V_2=s^2V,\qquad P_2=s^2P.
\]
Thus $P,V_2$ lie at level $d+2$ and $P_2$ lies at level $d+4$.
Using two tags at level $d+2$, form the Kuratowski-tagged coproduct
\[
Q=(P\times_K\{0\})\cup(V_2\times_K\{1\}),
\]
also at level $d+4$.  Even and odd elements of the typed Peano carrier
give a surjection $P_2\surj Q$.  In coordinates, for
$p=\langle\!\langle u,m\rangle\!\rangle\in P$, send
$s^2\langle\!\langle u,2m\rangle\!\rangle$ to the left-tagged copy of
$p$, and send
$s^2\langle\!\langle u,2m+1\rangle\!\rangle$ to the right-tagged copy of
$s^2g(u)$.  Surjectivity of $g$ makes the odd branch cover $V_2$.
Left tagging gives $P_2\inj Q$.  WPP therefore gives $Q\inj P_2$.

Restrict to the right summand.  Pairing naturality identifies it with
$s^2V_2=s^4V$, so the restriction is $s^4V\inj s^2P$.  Both endpoints
carry a common twofold support.  Lemma~\ref{lem:local-lowering} lowers the
restricted graph twice and gives \eqref{eq:typed-padding}.  The arbitrary
WPP witness is lowered only after both endpoints have been restricted to
named supports.
\end{proof}

\begin{lemma}[Typed Hartogs product]
\label{lem:typed-hartogs-product}
For displayed carriers $A,E$ at one level, suppose the larger of their
typed Hartogs orders is above the typed Peano order.  Then the ordinary
projection argument gives
\begin{equation}
 \mathbf H(A\times_KE)
 \inj
 \max_{\rm wo}\{s^2\mathbf H(A),s^2\mathbf H(E)\}.
 \label{eq:typed-hartogs-product}
\end{equation}
In particular, if the typed Hartogs order of $X$ is above the typed Peano
order and
\[
 P=(X\times_KX)\times_Ks^2\Omega,
\]
then
\begin{equation}
 \mathbf H(P)\inj s^4\mathbf H(X).
 \label{eq:typed-hartogs-special}
\end{equation}
\end{lemma}

\begin{proof}
Let $W$ well-order a subset $S$ of $A\times_KE$.  The supported coordinate
maps send $\langle\!\langle a,e\rangle\!\rangle$ to $s^2a$ and $s^2e$.
Order each coordinate image by $W$-least preimage.  No choice is used.
The coordinate map embeds $S$ into the type-level product of the two
well-ordered images.  If at least one image is infinite, the product of
their well-ordered cardinals is their maximum.  If both images are finite,
their product is finite and hence lies below the assumed infinite larger
Hartogs order.  Thus the transported code lies below the larger of
$\mathbf H(s^2A)$ and $\mathbf H(s^2E)$,
which are $s^2\mathbf H(A)$ and $s^2\mathbf H(E)$ by
\eqref{eq:hartogs-support-naturality}.

This proves pointwise coverage of source Hartogs codes.  Put the product
base and the selected supported base into their union and apply
Lemma~\ref{lem:common-ambient-quotient-graph}; its canonical graph proves
\eqref{eq:typed-hartogs-product} without selecting all pointwise maps at
once.  Apply the result first to $X,X$ and then to
$X\times_KX,s^2\Omega$.  The diagonal
$s^2X\inj X\times_KX$ shows that the second application also satisfies
the infinitude hypothesis.  Raising the first product comparison twice
gives
\[
 s^2\mathbf H(X\times_KX)\inj s^4\mathbf H(X).
\]
Moreover, the threshold assumption says that the typed Peano order lies
below $\mathbf H(X)$; equivalently, the typed Peano carrier injects into
$X$, and Hartogs monotonicity gives
$\mathbf H(\Omega)\inj\mathbf H(X)$.  Support naturality identifies
$s^2\mathbf H(s^2\Omega)$ with $s^4\mathbf H(\Omega)$, which therefore
embeds into $s^4\mathbf H(X)$.  Thus both coordinates of
the second maximum embed into $s^4\mathbf H(X)$, proving
\eqref{eq:typed-hartogs-special}.
\end{proof}

\begin{lemma}[Supported first-difference envelope]
\label{lem:supported-first-difference-envelope}
Assume WPP and that the typed Hartogs order of $X$ is above the typed
Peano order.  If $D$ is well ordered and
\[
 X\times_KX\surj sD,
\]
then
\begin{equation}
 sD\inj\mathbf H(X).
 \label{eq:supported-fd-envelope}
\end{equation}
\end{lemma}

\begin{proof}
Lemma~\ref{lem:typed-wpp-padding}, with
$U=X\times_KX$, $V=sD$, and $\Omega_2=s^2\Omega$, gives
\[
 s^3D\inj P=(X\times_KX)\times_Ks^2\Omega.
\]
Transport the well-order of $s^3D$ to its range.  The native typed
Hartogs initial-segment map has domain $s^5D=s^4(sD)$, and therefore gives
\[
 s^5D\inj\mathbf H(P)\inj s^4\mathbf H(X),
\]
where the second map is \eqref{eq:typed-hartogs-special}.  The two
endpoints have a common fourfold support.  Lowering them together proves
\eqref{eq:supported-fd-envelope}.  No iterate is removed from only one
side.
\end{proof}

\section{The bounded first-difference calculation}
\label{sec:bounded-calculation}

The calculation in this section records the ordinary carrier mathematics
which the explicit typed maps of Section~3 realize.  No collapse to von
Neumann ordinals and no simultaneous choice of witnesses is used.

\begin{lemma}[Countable padding]\label{lem:padding}
Assume WPP.  If $U\surj V$, then
\begin{equation}
 V\inj U\times\omega.
 \label{eq:padding}
\end{equation}
\end{lemma}

\begin{proof}
The empty case is immediate.  Otherwise, using
\(\omega\cong(\omega\times\{0\})\sqcup(\omega\times\{1\})\), construct a
surjection
\[
 U\times\omega\surj (U\times\omega)\sqcup V.
\]
The left summand injects into this disjoint union.  Apply
\eqref{eq:wpp-local}; the resulting injection restricts to $V$.
\end{proof}

\begin{lemma}[Carrier first-difference compression]
\label{lem:first-difference}
Let $(\mathcal A,\prec)$ be a well-ordered set of distinct subsets of
$X$.  There is a well-ordered set $D\subseteq\mathcal A$ and an injection
\begin{equation}
 \mathcal A\inj\Pow^2D.
 \label{eq:first-difference}
\end{equation}
If $D\neq\varnothing$, there is also a surjection $X^2\surj D$.  If
$D=\varnothing$, then $\mathcal A\subseteq\{\varnothing,X\}$.  Every
set and graph in the construction is obtained by bounded separation from
a displayed finite product or power set.
\end{lemma}

\begin{proof}
For $x\in X$, put
\[
 r_x=\{A\in\mathcal A:x\in A\}.
\]
Let
\[
 D=\{\text{the $\prec$-least member of }
       r_x\mathbin\triangle r_y: r_x\neq r_y\}
 \subseteq\mathcal A.
\]
The set $D$ inherits the well-order $\prec$.  If $D\neq\varnothing$,
send $(x,y)$ to the least difference when $r_x\neq r_y$, and to the
fixed $\prec$-least member of $D$ otherwise.  This is a surjection
$X^2\surj D$.

Put
\[
 q_x=r_x\cap D.
\]
First differences ensure
\begin{equation}
 q_x=q_y\quad\Longrightarrow\quad r_x=r_y.
 \label{eq:row-separation}
\end{equation}
For $A\in\mathcal A$, define
\[
 C_A=\{q_x:x\in A\}\in\Pow^2D.
\]
The map $A\mapsto C_A$ is injective.  Indeed, equality of two columns and
\eqref{eq:row-separation} first gives inclusion of the corresponding
members of $\mathcal A$ and then, by symmetry, equality.  If $D$ is
empty, all rows are equal, so each member of $\mathcal A$ is either
$\varnothing$ or $X$; the same column map still embeds $\mathcal A$ into
$\Pow^2D$.  No transitive collapse, Replacement, or simultaneous
selection of witnesses is used.
\end{proof}

\begin{lemma}[Hartogs products]\label{lem:products}
If $U,V$ are infinite, then
\[
 \aleph(U\times V)=\max\{\aleph(U),\aleph(V)\}.
\]
In particular, if $\aleph(X)>\omega$, then
\begin{equation}
 \aleph(X^2\times\omega)=\aleph(X).
 \label{eq:product-special}
\end{equation}
\end{lemma}

\begin{proof}
The lower bound follows by fixing an element in each factor.  Conversely,
if an ordinal $\gamma$ injects into $U\times V$, order each projection
of the range by least preimage.  Their order types are some
\(\alpha<\aleph(U)\) and $\beta<\aleph(V)$, and
\(\gamma\inj\alpha\times\beta\).  Multiplication of infinite
well-ordered cardinals gives the upper bound.  If
$\aleph(X)>\omega$, then $X$ is infinite and, since Hartogs numbers are
initial ordinals, the countable factor contributes no larger Hartogs
number.  Applying the result twice gives \eqref{eq:product-special}.
\end{proof}

\begin{proposition}[The single stabilized Hartogs injection]
\label{prop:stabilized}
Assume WPP.  If $\aleph(X)>\omega$, then
\begin{equation}
 \aleph(\Pow X)
 \leq
 \aleph\bigl(\Pow^2(\aleph(X))\bigr).
 \label{eq:stabilized}
\end{equation}
Moreover, the witnessing injection uses only the finite carriers and maps
displayed in Sections~3 and~4.
\end{proposition}

\begin{proof}
Both sides of \eqref{eq:stabilized} are ordinals, so they are linearly
comparable.  Suppose the desired injection fails.  Then
\[
 \lambda:=\aleph\bigl(\Pow^2(\aleph(X))\bigr)
 <\aleph(\Pow X),
\]
and hence there is an injection $j:\lambda\inj\Pow X$.  Well-order the
range $\mathcal A=\im(j)$ by least preimage and apply
Lemma~\ref{lem:first-difference}.  It gives a well-ordered carrier $D$ and
an injection $\lambda\inj\mathcal A\inj\Pow^2D$.

If $D\neq\varnothing$, Lemma~\ref{lem:padding} applied to the displayed
surjection $X^2\surj D$ gives $D\inj X^2\times\omega$.
Lemma~\ref{lem:products} and $\aleph(X)>\omega$ give
\[
 \aleph(X^2\times\omega)=\aleph(X).
\]
Because $D$ is well ordered, $D\inj\aleph(X)$.  The same conclusion is
immediate when $D=\varnothing$.  Double direct image now
gives
\[
 \lambda\inj\Pow^2D
       \inj\Pow^2\bigl(\aleph(X)\bigr),
\]
contrary to the definition of $\lambda$.  This proves
\eqref{eq:stabilized}.  The proof fixes only the one ordinal $\lambda$;
it never collects a family of first-difference witnesses.
\end{proof}

\section{The common-block pullback to NF}
\label{sec:pullback}

\begin{theorem}[Unconditional fixed-block replay]
\label{thm:unconditional-fixed-block-replay}
Assume NF+WPP, and suppose the typed Hartogs order of $X$ is above the
typed Peano order.  Then there is a set-coded injection
\begin{equation}
 s^3\mathbf H(\Pow X)
 \inj
 \mathbf H\bigl(\Pow^2(\mathbf H(X))\bigr).
 \label{eq:common-block-injection}
\end{equation}
\end{theorem}

\begin{proof}
Put
\[
 A_0=s^3(\Pow X),\qquad
 B=\Pow^2(\mathbf H(X)),\qquad C=A_0\cup B.
\]
All three bases lie at relative level $4$, while
$\mathbf H(A_0)$, $\mathbf H(B)$, and $\mathbf H(C)$ lie at level $6$.
Consider the set relation
\begin{equation}
 F=(\iota_B^C)^{-1}\circ\iota_{A_0}^C\circ\theta_{\Pow X}^3.
 \label{eq:canonical-final-graph}
\end{equation}
We prove that it is total; its definition already contains no choices.

Fix $x\in s^3\mathbf H(\Pow X)$.  Reverse coverage in
\eqref{eq:hartogs-support-naturality} supplies a well-order code on a
supported carrier $s^3I\subseteq A_0$, where $I\subseteq\Pow X$.
Lower its relation along the support isomorphism.  Apply the formulas in
Lemma~\ref{lem:first-difference} with the levels in
Table~\ref{tab:intrinsic-certificate}.  They give a well-ordered
$D\subseteq I$ and the typed maps
\begin{equation}
 X\times_KX\surj sD,
 \qquad
 s^2I\inj\Pow^2D,
 \label{eq:typed-first-difference-maps}
\end{equation}
when $D\neq\varnothing$; when $D=\varnothing$ the second map remains and
the first is unnecessary.

Lemma~\ref{lem:supported-first-difference-envelope}, or the empty map in
the empty case, gives $k:sD\inj\mathbf H(X)$.  Raise the column map once,
use power naturality, and take two direct images under $k$:
\begin{equation}
 s^3I
 \inj s(\Pow^2D)
 \cong\Pow^2(sD)
 \inj\Pow^2(\mathbf H(X))=B.
 \label{eq:typed-column-composite}
\end{equation}
Push the supported well-order forward along this injection to its range in
$B$, and let $y\in\mathbf H(B)$ be the resulting quotient class.  In the
common ambient base $C$, the original code and the transported code are
isomorphic.  Hence
\[
 \iota_{A_0}^C\bigl(\theta_{\Pow X}^3(x)\bigr)=\iota_B^C(y).
\]
This proves
\(\iota_{A_0}^C``\mathbf H(A_0)
  \subseteq\iota_B^C``\mathbf H(B)\)
pointwise.  Lemma~\ref{lem:common-ambient-quotient-graph} now says that
\eqref{eq:canonical-final-graph} is a total injection with the endpoints
in \eqref{eq:common-block-injection}.

The first-difference injection, the padding witness, and the transported
code were used only to prove range coverage.  None occurs in the definition
of $F$, so no family of witnesses is chosen.
\end{proof}

\begin{lemma}[Homogeneous stabilized inequality]\label{lem:homogeneous}
Assume NF together with WPP.  Let $A$ be a typed set, write $a=|A|$, and
suppose the corresponding Hartogs number is above the typed natural-number
cardinal.  Then
\begin{equation}
 h(ea)\leq h\bigl(e^2(ha)\bigr).
 \label{eq:homogeneous}
\end{equation}
\end{lemma}

\begin{proof}
Apply Theorem~\ref{thm:unconditional-fixed-block-replay}.  The endpoint
objects and their typed cardinals are as follows.
\begin{table}[ht]
\centering
\caption{Typed representatives and relative depths for the stabilized injection.}
\label{tab:ledger}
\begin{tabular}{@{}lcl@{}}
\toprule
typed carrier & relative depth & cardinal \\
\midrule
$A$ & $0$ & $a$ \\
$\Pow A$ & $1$ & $T(ea)$ \\
$\mathbf H(A)$ & $2$ & $T^2(ha)$ \\
$\mathbf H(\Pow A)$ & $3$ & $T^3h(ea)$ \\
$\Pow(\mathbf H(A))$ & $3$ & $T^3e(ha)$ \\
$\Pow^2(\mathbf H(A))$ & $4$ & $T^4e^2(ha)$ \\
$s^3(\mathbf H(\Pow A))$ & $6$ & $T^6h(ea)$ \\
$\mathbf H(\Pow^2(\mathbf H(A)))$ & $6$ & $T^6h(e^2(ha))$ \\
\bottomrule
\end{tabular}
\end{table}

The last two rows are the aligned endpoints of the actual set graph
\eqref{eq:canonical-final-graph}.  If expansion into a chosen primitive
pair code adds $r$ levels, it adds the same $r$ to both endpoints and gives
\[
 T^{6+r}h(ea)\leq T^{6+r}h(e^2(ha)).
\]
Since $T$ reflects cardinal comparison, cancelling the common iterate
$T^{6+r}$ proves \eqref{eq:homogeneous}.  No comparison of the form
$TK\leq K$, and no inverse to $T$, is used.
\end{proof}

\begin{remark}
The common-level step is essential.  If one normalizes the source and
target separately and then drops a $T$ from only one side, one obtains a
stronger predecessor inequality which does not follow from the product
calculation.  Lemma~\ref{lem:homogeneous} avoids that error by lowering the
actual final injection graph only after its endpoints have been aligned.
\end{remark}

\section{The finite-orbit contradiction}

Let $U=|V|$, the cardinality of the universal set.  Since
\(\Pow V=V\), the homogeneous exponential satisfies
\begin{equation}
 e(T^{n+1}U)=T^nU
 \qquad(n\text{ a concrete natural number}).
 \label{eq:universal-chain}
\end{equation}
Cantor's theorem gives $T^{n+1}U<T^nU$.

Define, wherever the displayed terms exist,
\begin{equation}
 F(\beta)=h\bigl(h(e^2\beta)\bigr).
 \label{eq:def-F}
\end{equation}
This operation is homogeneous and nondecreasing.  For every well-orderable
\(\beta\) in its domain,
\begin{equation}
 \beta<h(e^2\beta)<h\bigl(h(e^2\beta)\bigr)=F(\beta),
 \label{eq:F-strict}
\end{equation}
so $F(\beta)>\beta$.  Indeed, the double-singleton map gives
$\beta\leq e^2\beta$; because $\beta$ is well orderable, the defining
Hartogs property gives the first strict inequality.  The inner Hartogs
value is itself well orderable, giving the second.

\begin{lemma}[Domain width]\label{lem:domain-width}
There is a concrete natural number $w$ such that
\begin{equation}
 \operatorname{ran}(T^w)\cap\mathrm{WCard}
 \subseteq\operatorname{dom}F,
 \label{eq:F-domain-width}
\end{equation}
and $F$ commutes with $T$ on that domain.
\end{lemma}

\begin{proof}
Choose the fixed typed definitions of the two power-set and two Hartogs
constructions in \eqref{eq:def-F}, and let $w$ exceed the largest relative
type displacement among their input and output carriers.  A well-orderable
cardinal in $\operatorname{ran}(T^w)$ then has a representative at the
input level of this fixed diagram.  Power-set naturality and the two-way
Hartogs-code naturality proved in Lemma~\ref{lem:local-lowering} construct
each successive value and prove the commutation law.  This is a structural
induction on the four displayed constructors, not an inverse use of $T$.
\end{proof}

We first isolate the orbit argument.

\begin{lemma}[Least-orbit obstruction]\label{lem:orbit}
Let $F$ be a homogeneous nondecreasing partial operation on
well-orderable cardinals.  Suppose $F$ is defined on the initial segment
through $\alpha$, $F(\beta)>\beta$ there, and
\begin{equation}
 T\alpha<\alpha<F(T\alpha).
 \label{eq:sandwich}
\end{equation}
If no Cantorian well-orderable cardinal $\beta=T\beta$ satisfies
$\beta\geq\alpha$, then NF proves a contradiction.
\end{lemma}

\begin{proof}
For $a\leq\alpha$, let $\Gamma(a)$ be the least well-orderable
cardinal $\delta\leq a$ whose finite forward $F$-orbit reaches $a$.
The candidate set is nonempty because $a$ qualifies at time zero.
Candidate transport is two-sided.  Shifting a start and its hitting time
gives a candidate for $Ta$.  Conversely, if $\delta\leq Ta$ is a candidate
for $Ta$, initiality of $\operatorname{ran}T$ supplies
$\delta=T\epsilon$ with $\epsilon\leq a$.  If $m$ is a hitting-time
witness, surjectivity of $T$ on the naturals gives $m=Tn$.  Homogeneity and
order reflection then give
\[
 T(F^n(\epsilon))=F^{Tn}(T\epsilon)=F^m(\delta)\geq Ta,
 \qquad F^n(\epsilon)\geq a.
\]
Thus $\epsilon$ is a candidate for $a$.  Leastness in both directions
gives, whenever $a\leq\alpha$ and $Ta\leq\alpha$,
\begin{equation}
 T\Gamma(a)=\Gamma(Ta).
 \label{eq:gamma-shift}
\end{equation}

An orbit reaches $T\alpha$ if it reaches $\alpha$, because
$T\alpha<\alpha$.  Conversely, after first reaching a value
$y\geq T\alpha$ below $\alpha$, its next value satisfies
\[
 F(y)\geq F(T\alpha)>\alpha.
\]
Thus reaching the two thresholds is equivalent.  Also $T\alpha$ reaches
$\alpha$ in one step, so the least start for $\alpha$ is at most
$T\alpha$.  The different upper bounds in the two definitions therefore
do not change the least candidate:
\[
 \Gamma(T\alpha)=\Gamma(\alpha).
\]
Put $\gamma=\Gamma(\alpha)$.  Equation~\eqref{eq:gamma-shift} gives
\begin{equation}
 T\gamma=\gamma.
 \label{eq:gamma-fixed}
\end{equation}

Let $m$ be the least natural number such that
\(F^m(\gamma)\geq\alpha\), and let $k$ be the least time at which the
same orbit reaches $T\alpha$.  We claim $k=Tm$.  Homogeneity and
\eqref{eq:gamma-fixed} show that $Tm$ is a lower-threshold hit.  If
$j<Tm$ were an earlier one, write $j=Tq$ by surjectivity of $T$ on the
naturals.  Order reflection gives $q<m$, whereas
\[
 T(F^q(\gamma))=F^{Tq}(T\gamma)=F^j(\gamma)\geq T\alpha
\]
would imply $F^q(\gamma)\geq\alpha$, contradicting minimality of $m$.
Thus $k=Tm$.

Since the upper threshold is larger, $k\leq m$.  At time $k$ the orbit is
at least $T\alpha$; if it is still below $\alpha$, one further application
of $F$ passes $\alpha$, and if not, none is needed.  Hence
\[
 Tm=k\leq m\leq k+1=Tm+1.
\]
Discreteness of the natural-number order gives
$m=Tm$ or $m=Tm+1$.

In the first case, $\beta=F^m(\gamma)$ satisfies
\[
 T\beta=F^{Tm}(T\gamma)=F^m(\gamma)=\beta
\]
and $\beta\geq\alpha$, contradicting the hypothesis about Cantorian
well-orderable cardinals.  The second case is impossible because $m$ and
$Tm$ have the same parity.  The argument does not require $m$ to be a
standard natural number.
\end{proof}

\begin{lemma}[Universal-chain threshold]\label{lem:threshold}
Assume NF together with WPP.  For a sufficiently large concrete $n$,
the cardinal
\begin{equation}
 \alpha=h(T^nU)
 \label{eq:def-alpha}
\end{equation}
satisfies all hypotheses of Lemma~\ref{lem:orbit} for the operation
\eqref{eq:def-F}.
\end{lemma}

\begin{proof}
We first justify the domain.  Let $w$ be supplied by
Lemma~\ref{lem:domain-width}, and put
\[
 D=|\mathbf H(V)|=h(T^2U),
\]
which is the cardinal of an actual typed Hartogs object.  For $n\geq2$,
shift coherence gives
\begin{equation}
 h(T^nU)=T^{n-2}D.
 \label{eq:alpha-depth}
\end{equation}
Choose $n$ so large that $n-2\geq w$.  Then
$\alpha=T^{n-2}D$ lies in $\operatorname{ran}(T^w)$; initiality of that
range on well-orderable cardinals and \eqref{eq:F-domain-width} give
\[
 [0,\alpha]\cap\mathrm{WCard}
 \subseteq\operatorname{ran}(T^w)
 \subseteq\operatorname{dom}F.
\]
This replaces the meaningless shorthand $h(U)$ sometimes used in informal
versions of the argument.

Take a typed representative $A$ of cardinal $T^{n+1}U$.  At this
fixed concrete height a typed copy of the natural numbers injects into
$A$, so the Hartogs number corresponding to $A$ is above the natural
cardinal.  Lemma~\ref{lem:homogeneous} applies with $a=T^{n+1}U$.
Using \eqref{eq:universal-chain} and commutation with $T$,
\[
 \alpha
 =h(e(T^{n+1}U))
 \leq h(e^2(h(T^{n+1}U)))
 =h(e^2(T\alpha))
 <F(T\alpha).
\]

Monotonicity and $T^{n+1}U<T^nU$ give
\[
 T\alpha=h(T^{n+1}U)\leq h(T^nU)=\alpha.
\]
If a well-orderable cardinal $\beta$ is Cantorian, then
\[
 \beta=T^n\beta\leq T^nU,
\]
and the defining property of $h(T^nU)=\alpha$ yields
\(\beta<\alpha\).  Thus no Cantorian well-orderable cardinal is at least
$\alpha$; in particular, the preceding inequality is strict.  Hence
\[
 T\alpha<\alpha<F(T\alpha),
\]
and all hypotheses of Lemma~\ref{lem:orbit} hold.
\end{proof}

\begin{proof}[Proof of Theorem~\ref{thm:main}]
Assume NF+WPP.  The carrier first-difference construction and the finite
typed replay in Theorem~\ref{thm:unconditional-fixed-block-replay} give
Lemma~\ref{lem:homogeneous}.  The universal-chain specialization in
Lemma~\ref{lem:threshold} then meets the hypotheses of
Lemma~\ref{lem:orbit}, a contradiction.  Therefore
\(\NF\vdash\neg\WPP\).
\end{proof}

\begin{corollary}
NF refutes every choice principle that implies WPP; in particular, it
refutes the usual partition principle and the full axiom of choice.
\end{corollary}

\begin{proof}
Immediate from Theorem~\ref{thm:main} and the stated implications.
\end{proof}

\section{What the proof uses}

The reduction uses an ordinary bounded carrier calculation only as a
finite derivation template.  Its typed content consists of:
\begin{enumerate}[(1)]
\item one countable-padding graph;
\item one well-ordered first-difference carrier and one column injection;
\item finitely many power-set monotonicity maps;
\item one Hartogs-product comparison;
\item one final injection between two Hartogs quotient objects.
\end{enumerate}
Lemma~\ref{lem:local-lowering} realizes each local support-lowering step;
Lemmas~\ref{lem:typed-wpp-padding}--\ref{lem:supported-first-difference-envelope}
assemble the padding and product maps; and the common-ambient graph
\eqref{eq:canonical-final-graph} supplies the aligned final injection.
Thus the proof requires neither coded formulas and satisfaction nor a
direct-limit membership relation.  It is a finite collection of ordinary
stratified NF constructions, with every temporary witness eliminated
before the final graph is formed.

This economy also makes the type ledger auditable.  The only cancellation
in the homogeneous pullback is
\[
 T^{6+r}h(ea)\leq T^{6+r}h(e^2ha)
 \quad\Longrightarrow\quad
 h(ea)\leq h(e^2ha),
\]
with the same exponent on both sides.  This is the mathematical reason for
using the explicit finite-block route rather than the shorter direct
product argument.

\section*{Acknowledgments}
I thank M. Randall Holmes for the finite-orbit formulation of the final
argument and for discussions of typed realizations of NF.

\begin{thebibliography}{99}

\bibitem{Forster}
Thomas E. Forster,
\emph{Set Theory with a Universal Set}, second edition,
Oxford Logic Guides, Oxford University Press, 1995.

\bibitem{HolmesTangled}
M. Randall Holmes,
The equivalence of NF-style set theories with ``tangled'' type theories;
the construction of omega-models of predicative NF (and more),
\emph{Journal of Symbolic Logic} \textbf{60} (1995), 178--189.

\bibitem{HolmesBook}
M. Randall Holmes,
\emph{Elementary Set Theory with a Universal Set},
revised online edition, 2012.

\bibitem{HolmesWilshaw}
M. Randall Holmes and Sky Wilshaw,
\emph{NF is Consistent}, arXiv:1503.01406v23, 2025.

\bibitem{Specker}
E. P. Specker,
The axiom of choice in Quine's ``New Foundations for Mathematical Logic'',
\emph{Proceedings of the National Academy of Sciences USA}
\textbf{39} (1953), 972--975.

\end{thebibliography}

\end{document}
